\section{Runtime Design}\label{sec:design}
\subsection{Facts, projections, and physical activation}
Figure~\ref{fig:architecture} shows the separation used in our analysis. Domain inputs and observed outcomes contribute durable facts. Reducers interpret those facts with lifecycle policy, ownership, and time. The driver requests scans and arranges activation; a scan is not itself a new task definition. Agent turns produce real tool requests and observations, and delivery decisions refer back to identifiable evidence.

\begin{figure}[htbp]\centering
\begin{tikzpicture}[font=\small,>=Latex,box/.style={draw,rounded corners=2pt,align=center,minimum height=12mm,text width=37mm},node distance=8mm]
\node[box] (facts) {Durable facts\\and lifecycle revisions};
\node[box,right=of facts] (reduce) {Projection\\facts, policy, leases, time};
\node[box,right=of reduce] (driver) {Driver and activation\\hints, scans, wakes};
\node[box,below=of facts] (evidence) {Exact evidence\\read, select, publish};
\node[box,below=of reduce] (turn) {Agent turn\\tool request / observation};
\node[box,below=of driver] (world) {External environment\\operation-specific effects};
\draw[->] (facts)--(reduce);\draw[->] (reduce)--(driver);
\draw[->] (driver)--(turn);\draw[->] (turn)--(world);
\draw[->] (turn)--(evidence);\draw[->] (evidence)--(facts);
\draw[->] (world.south) -- ++(0,-5mm) -| (turn.south);
\node[align=center,font=\footnotesize,below=12mm of turn] {Delivery decision references evidence and accepted requirement revisions.};
\end{tikzpicture}

\caption{Analytical view of selected \oc{} paths. Arrows denote information or authority dependencies, not a claim of a single atomic pipeline. External effects cross the database boundary.}\label{fig:architecture}
\end{figure}

The root-ingress reducer illustrates the distinction between a pending request and permission to continue it. In the inspected source, operator abandonment and integrity faults are classified explicitly; lifecycle epochs can make older ingress inapplicable; validated decisions can resolve it; and live leases identify an already-owned activation. These alternatives matter because blindly re-running a queued instruction after a restart can apply it to a different lifecycle interval. The reducer and wake classifier expose whether progress depends on a deadline, an infrastructure fact, or operator input. A host fault is not evidence that a domain task was successfully completed.

The driver addresses a separate problem: durable work may exist without a surviving in-memory notification. When a request arrives, the local task entry's revision increases. An existing scan owner absorbs the notification; otherwise the driver attempts capacity admission. The owner scans for a bounded number of passes, records requested wake instants, and rechecks when revisions change. Faults and non-progress are paced with backoff. The periodic live-task enumeration retains rejected admissions in its current page rather than committing past them. This mechanism supplies rediscovery opportunities without making local notification counts the business state.

The schematic procedure below is an abstraction of these selected paths. It deliberately omits policy-specific reducer branches, retry coefficients, and project entry-point wiring. It is not an executable replacement for the runtime.
\begin{quote}\small
\begin{enumerate}
\item On a hint for task $\tau$, increment its local revision. Coalesce with a running owner or request a scan slot.
\item Under an admitted owner, read the current fact revision and invoke the scan against durable state.
\item Collect activation and wake information. If new input arrived, rescan within the pass budget; otherwise stop the pass loop.
\item On a bounded failure or lack of progress, retain the condition and arrange a later attempt under backoff.
\item Release local scan ownership and arm the earliest retained wake, if any.
\item Independently enumerate live tasks. Keep capacity-rejected entries eligible for later admission.
\end{enumerate}
\end{quote}

Local scan ownership and database activation leases serve different purposes. The former avoids redundant work within a driver instance; the latter guards selected durable transitions across physical contenders. Neither substitutes for project-scoped data access or for exact execution lineage. A complete deployment audit would need to trace every production entry point and ownership boundary; this paper does not claim to have formally verified that entire surface.

\subsection{Ownership checked at the durable transition}
The lease module represents a target, an owner occurrence, a lease identity, activation time, and expiry. Acquisition reads the current lease in a transaction and normally refuses a new contender while that lease is unexpired. Explicit supersession is a separate authority-bearing case. Renewal requires the current identity and owner, updates the expected expiry, and appends a grant record. Release expires the exact current lease. Thus the durable structure includes both mutable expiry and grant history.

A stale worker is rejected only where a protected transition actually rechecks ownership. For example, the dispatch-lineage commitment path accepts an admission owner, asserts its lease when present, requires a durable descriptor for the exact dispatch turn, and consumes the admission in the same transaction. The optional admission parameter is material: this code inspection does not imply that every invocation carries an admission lease. Likewise, task-completion closure uses a lifecycle-scoped target and translates failed ownership assertions into an explicit closure conflict.

Figure~\ref{fig:recovery} separates database rejection from an external effect's outcome. If a worker performs an external operation before it loses its lease, a later database rejection cannot erase that operation. If an acknowledgement is lost, the runtime may know that it lacks a result while the external service has already committed the effect. Safe recovery requires an operation-specific reconciliation or idempotency contract; the existence of a lease is insufficient.

\begin{figure}[htbp]\centering
\begin{tikzpicture}[font=\small,>=Latex]
\node (old) at (0,0) {Owner A};\node (db) at (4.3,0) {Database};\node (ext) at (8.7,0) {External service};\node (new) at (13,0) {Owner B};
\foreach \n in {old,db,ext,new} {\draw[densely dashed,gray] (\n.south)--++(0,-48mm);}
\draw[->] (0,-.7)--node[above,font=\footnotesize]{acquire lease A}(4.3,-.7);
\draw[->] (0,-1.5)--node[above,font=\footnotesize]{perform effect}(8.7,-1.5);
\node[font=\footnotesize,align=center] at (8.7,-2.15) {acknowledgement lost};
\draw[->] (13,-2.8)--node[above,font=\footnotesize]{acquire after expiry}(4.3,-2.8);
\draw[->] (0,-3.5)--node[above,font=\footnotesize]{assert stale lease A}(4.3,-3.5);
\draw[->] (4.3,-4.15)--node[above,font=\footnotesize]{fence conflict}(0,-4.15);
\node[font=\footnotesize,align=center] at (10.6,-4.1) {Effect outcome needs\\independent reconciliation.};
\end{tikzpicture}

\caption{Schematic failure sequence, not an observed experimental trace. The stale owner's guarded database write is rejected; whether its earlier external effect occurred requires separate evidence.}\label{fig:recovery}
\end{figure}

\subsection{Evidence is read, selected, and published}
The artifact catalog provides exact references to engine artifact revisions, task snapshots, and resources. The read path checks identity and integrity rather than interpreting an unversioned name as the latest content. Text may be read in bounded windows. The recorded windows must collectively cover the exact reference before it qualifies as a complete read in the relevant scope. A complete transport record means that the bytes were delivered through that interface; it does not reveal how a model attended to them.

The publication function accepts observed and selected locator sets from its host caller. Every declared source locator must belong to both sets. Resource checks also relate the publication to its task and validate referenced resource content. The two inspected public call paths supply these sets differently: the native tool derives scoped read and selection facts preceding the current publication, whereas the plugin host tracks reads and selections within the invocation. Both call the same catalog publication function. The authority of the input sets depends on those collectors; the function cannot prove their origin merely from a supplied array.

Figure~\ref{fig:evidence} illustrates the resulting relation. Exact references prevent a reviewer from silently being credited with a later revision. Explicit selection records a producer's declared use. Yet neither guarantees comprehensive citation: a producer can omit a relevant source or misunderstand one. Domain review and independent acceptance remain separate activities.

\begin{figure}[htbp]\centering
\begin{tikzpicture}[font=\small,>=Latex,box/.style={draw,rounded corners=2pt,align=center,text width=24mm,minimum height=14mm},node distance=5mm]
\node[box] (ref) {Exact source\\identity, revision, digest};
\node[box,right=of ref] (read) {Complete read\\byte coverage};
\node[box,right=of read] (select) {Selection\\declared use};
\node[box,right=of select] (pub) {Publication\\source membership};
\node[box,right=of pub] (accept) {Decision\\accepted revisions};
\draw[->] (ref)--(read);\draw[->] (read)--(select);\draw[->] (select)--(pub);\draw[->] (pub)--(accept);
\node[draw,align=center,below=9mm of select,text width=75mm] (oracle) {Independent oracle: correctness, required evidence, and task acceptance};
\draw[->] (pub.south) -- ++(0,-4mm) -| (oracle.north);
\end{tikzpicture}

\caption{Evidence path for one declared source. Byte coverage, declared use, publication, and acceptance have different semantics. The independent oracle evaluates the deliverable separately.}\label{fig:evidence}
\end{figure}

\subsection{Delivery decisions and responsibility}
The inspected completion-decision schema records orchestrator session/message identity, tool call/part identity, evidence locators, deliverable locators, accepted slice revision identifiers, workflow binding, and recording time. It includes host-derived references to outputs from terminal workflow nodes. A reader locates the decision for an exact task and terminal time. This makes an acceptance decision inspectable as a particular act with referenced material.

The closure establishes a recorded set and attribution, not a proof that the model judged correctly. An output from a terminal workflow node does not certify that all preceding nodes executed in order. A passing host checker certifies only the tested contract and observed execution context. Visual or domain judgments need their own evidence; tests over unrelated code or a final status label cannot replace them. This boundary permits meaningful accounting when the runtime records completion but an independent evaluator rejects the deliverable.

\section{Implementation and Audit Scope}\label{sec:implementation}
The source inspection is fixed to commit \code{bd2d6bd18eda05f35326adc231289e56d0a2dfcf}. The manuscript is stored in a later documentation-only state. Product source is not modified for this paper. Table~\ref{tab:implementation} provides the principal source anchors; the accompanying source map records symbols, line spans, and pre-existing claim identifiers. The symbols are preferable to floating line numbers when a reader inspects another revision.

\begin{table}[htbp]\centering\small
\begin{tabularx}{\linewidth}{@{}YY@{}}
\toprule
Module under \code{packages/opencorvus/src/} & Inspected responsibility \\
\midrule
\path{engine/task-root-ingress-reducer.ts} & Classify ingress facts and wake conditions. \\
\path{engine/task-control-driver.ts} & Coalesce hints, admit scans, retain wakes, enumerate live tasks. \\
\path{engine/control-lease.ts} & Acquire, assert, renew, and release exact ownership. \\
\path{engine/dispatch-lineage.ts} & Bind dispatch to a durable turn descriptor; guarded admission consumption. \\
\path{engine/task-completion-closure.ts} & Lifecycle-scoped ownership of completion closure. \\
\path{engine/completion-decision-facts.ts} & Typed acceptance references and terminal-time lookup. \\
\path{artifact-catalog/index.ts} & Exact reads, resource integrity, source membership at publication. \\
\path{agent/artifact-provenance-facts.ts} & Scoped complete reads and selected source facts. \\
\path{tool/artifact-catalog.ts}; \path{tool/plugin-tool-host.ts} & Supply provenance from actual tool/action scope. \\
\bottomrule
\end{tabularx}
\caption{Inspected implementation anchors. This table reports source responsibilities, not runtime benchmark validation.}\label{tab:implementation}
\end{table}

The repository uses TypeScript and database access through its storage abstraction; the lease acquisition wrapper uses an immediate transaction. These are implementation choices rather than prerequisites of the semantic model. The model requires an appropriate transactional boundary and persistent identity, not a particular programming language. Existing architecture documents describe capabilities and package projection, but this manuscript does not count packages, agents, providers, or tools as research contributions. Live provider projection, all production entry points, and full cross-project isolation have not been exercised for this study.

We preserve the repository's earlier evidence package as historical input. Its machine claim matrix remains the original claim authority; the new source map relates paper passages to that package and to inspected source. It does not upgrade old implementation statuses or acceptance receipts. Source inspection can establish that a check exists in a selected path. Establishing that all actual runs traverse the intended checks requires additional execution and coverage evidence.
